\section{결론}
\label{sec:conclusion}

추론 보강형 E2E 주행 데이터는 행동 이유를 제공하지만, 개별 CoC만 검사하면 이전 사건에서 시작된 지속 의무와 해제 상태를 보존할 수 없다. 본 논문은 연속 CoC를 행동 의도, 지속 의무, 해제ㆍ충족 조건, 필수 행동 순서와 삼값 증거를 갖는 공통 중간표현으로 변환하고, 상태형 Python 검사기와 같은 전이 의미론의 \uppaal 관찰 모델로 실행하는 방법을 제시하였다. 중심 관찰 모델은 사건 순서와 의무 상태를 검사하며, 실제 시간 기한은 별도의 보조 반응 계약에서만 다룬다.

자차 궤적과 연결된 다중-CoC 90장면에서 309개 인접 창을 구성하고 선택한 27개 장면을 네 검토자가 판정했다. 합의상 텍스트 모순은 0/27이었지만 합의 전 텍스트 일관성 Fleiss $\kappa$는 0.027이었고 모든 장면에서 적어도 한 항목의 의견이 갈렸다. 따라서 이 자연 표본은 실제 모순의 부재나 자동 검출 효과를 뒷받침하지 않는다.

결정적 합성 중간표현 40쌍에서는 상태형 규칙이 변형 40/40을, 사건별 규칙이 20/40을 모순으로 판정했고 기준 사례의 모순 판정은 모두 0/40이었다. 그러나 생성기가 상태형 판정을 사후조건으로 강제하므로 이는 독립 성능 평가가 아니라 선언한 네 규칙 유형의 회귀 시험 커버리지이다. 같은 전이 규칙에서 생성한 일곱 UPPAAL 시험 사례와 Python의 판정 불일치는 0/7이었으며, 이는 제한된 구현 일치와 최초 위반 상태 추적을 지지한다.

보조 단일사건 분석의 $134=125+4+5$ 분류, 반복 정책ㆍ탐지기 민감도와 장면 경계 시험은 반응 시간ㆍ출처 계약의 조건 의존성을 보여 주지만 중심 순서ㆍ상태 계약의 자연 오류 검출 성능을 보강하지는 않는다. 본 결과는 자연 자료에서의 모순 발생 빈도, CoC의 일반적 충실도, 실제 차량 안전 또는 학습 성능 향상을 입증하지 않는다. 후속 연구에는 합의 근거가 보존된 독립 자연 모순ㆍ비모순 자료, 독립적으로 구현한 기준 검사기, 객체 상태와 궤적 증거, 그리고 폐루프 평가가 필요하다.
